% !TeX root = ../main.tex
\subsection{تمرین \secnum{۵.۴}، فعالیت از روی شکل موج}
\label{sec:assignment-power-waveform}

\textbf{خواسته.}
برای سیگنال شکل \secnum{۵.۳۴} کتاب با ساعت \qty{1}{\giga\hertz}،
\lr{Activity Factor} را تعیین کنید. شکل \ref{fig:assignment-power-waveform}
بازه و لبه‌های صعودی را نشان می‌دهد
\cite[تمرین \secnum{۵.۴}، صفحه ۲۰۹]{weste}.

\begin{figure}[htbp]
  \centering
  % Shared redraw of the exercise waveform for the notes and submission.
\begin{latin}
\begin{tikzpicture}[x=1.08cm,y=1.30cm]
  \draw[->] (0,-0.30) -- (10.75,-0.30) node[right] {\(t\,(\unit{\nano\second})\)};
  \node[left] at (-0.08,1.24) {\(v/V_{DD}\)};
  \node[left] at (-0.08,0) {\rl{۰}};
  \node[left] at (-0.08,1) {\rl{۱}};
  \foreach \x/\ticklabel in {0/۰,1/۱,2/۲,3/۳,4/۴,5/۵,6/۶,7/۷,8/۸,9/۹,10/۱۰} {
    \draw (\x,-0.25) -- (\x,-0.36) node[below] {\rl{\ticklabel}};
  }
  \draw[gray,densely dashed] (1,-0.24) -- (1,0.50);
  \draw[gray,densely dashed] (4,-0.24) -- (4,0.50);
  \draw[very thick] (0,0) -- (0.75,0) -- (1.25,1) -- (1.75,1)
    -- (2.25,0) -- (3.75,0) -- (4.25,1) -- (7.75,1) -- (8.25,0) -- (10,0);
  \fill (1,0.50) circle (1.5pt) (4,0.50) circle (1.5pt);
\end{tikzpicture}
\end{latin}

  \caption{بازرسم شکل \secnum{۵.۳۴} کتاب با محور ولتاژ نرمال‌شده.
    نقطه‌ها دو گذار صعودی در بازه مشاهده را مشخص می‌کنند \cite{weste}.}
  \label{fig:assignment-power-waveform}
\end{figure}

\textbf{حل.}
دوره ساعت \(T_{\mathrm{clk}}=1/f=\qty{1}{\nano\second}\) است.
در فاصله‌ی صفر تا \qty{10}{\nano\second}، ده دوره کامل داریم.
خروجی نزدیک \qty{1}{\nano\second} و \qty{4}{\nano\second} از
صفر به یک می‌رود. شمار گذار صعودی \(N_{\uparrow}=2\) و شمار
دوره ساعت \(N_{\mathrm{clk}}=10\) است. بر اساس تعریف کتاب، حاصل چنین است
\cite[بخش \secnum{۵.۱.۲}، صفحه ۱۸۴]{weste}.
\[
  \alpha=\frac{N_{\uparrow}}{N_{\mathrm{clk}}}=\frac{2}{10}=0.2.
\]
در قرارداد این فصل، هر گذار صعودی یک رویداد شارژ است و شارژها
در صورت کسر شمرده می‌شوند.
